% !PRINTED 427
\weberpage{427}
%<<CONT
the congruence is identical, i.e. that all coefficients of $f(x)$
must be divisible by $n$.

This theorem is, as one sees, entirely analogous to the algebraic theorem that an equation cannot have more roots
than its degree indicates. But the other theorem, that every
equation also actually has that many roots, cannot be carried
over. A congruence of degree $m$ can have fewer roots, indeed
none at all. All the more remarkable is a particular
congruence in which the number of roots always equals the
degree, on the basis of a theorem which is called \gesperrt{Fermat's theorem}, and which we formulate here as follows.

\gesperrt{2. The congruences}
\weq{5}{x^n - x \equiv 0, \quad x^{n-1} - 1 \equiv 0 \quad (\mathrm{mod}\ n)}
\gesperrt{have, if $n$ is a prime, as many roots
as their degree amounts to, namely $n$ and $n-1$.}

The two assertions are not essentially different, for
the first of the congruences (5) has all the roots of the second and
in addition the root 0, which does not satisfy the second. Likewise
the second has all the roots of the first, with the exception of the
root 0.

Now since for the modulus $n$ there are altogether only $n$ different
numbers, it is thus to be proved that for every integer
$\alpha$ the congruence holds
\weq{6}{\alpha^n \equiv \alpha \quad (\mathrm{mod}\ n).}
This congruence is correct for $\alpha = 0$ and $\alpha = 1$. We
therefore prove it again by complete induction, deriving the congruence
\weq{7}{(\alpha + 1)^n \equiv \alpha + 1 \quad (\mathrm{mod}\ n)}
from congruence (6), assumed to be correct. This follows from the binomial theorem,
if one observes that all binomial coefficients, with the exception of the first and the last, which are equal to 1, namely
\[
n, \quad \frac{n(n-1)}{1\,.\,2}, \quad \frac{n(n-1)(n-2)}{1\,.\,2\,.\,3}, \dots
\]
are integers divisible by $n$, because the numerator, but not
the denominator, is divisible by $n$. Accordingly
\[
(\alpha + 1)^n \equiv \alpha^n + 1 \quad (\mathrm{mod}\ n),
\]
hence the formula (7) follows from the formula (6).
